\chapter[Thảo luận, kết luận và hướng phát triển]{Thảo luận, kết luận\\và hướng phát triển}
\label{chap:discussion}

\section{Thảo luận các phát hiện chính}
Tổng hợp toàn diện chuỗi thực nghiệm ở Chương~\ref{chap:results}, khóa luận đúc kết ba bài học thiết kế then chốt cho bài toán học sâu sống còn đa phương thức trên dữ liệu y tế cỡ mẫu vừa và nhỏ:

Thứ nhất, \emph{hiệu quả nén tham số của phân rã CP cần được nhìn nhận đúng đắn theo từng phạm vi kiến trúc cụ thể}. Việc thay thế lớp chiếu tích ngoài đầy đủ bằng cấu trúc phân rã CP giúp cắt giảm tới \textbf{98.38\%} số lượng tham số tại đúng ba lớp song tuyến đầu tiên (từ 229696 xuống chỉ còn 3712 tham số) và đạt mức giảm 90.90\% khi tính chung cho cả khối MLP tương tác. Tuy nhiên, mức giảm này chỉ tương ứng với khoảng 0.67\% tổng số tham số của toàn bộ mô hình do bộ mã hóa ảnh ResNet-18 3D chiếm tỷ trọng hoàn toàn áp đảo. Phát hiện này làm sáng tỏ rằng: lợi ích lớn nhất của việc ứng dụng phân rã CP không nằm ở việc thu nhỏ kích thước lưu trữ của toàn mạng, mà tập trung ở việc \emph{ngăn chặn triệt để nguy cơ quá khớp cục bộ tại tầng hợp nhất đa phương thức} — vị trí vốn nhạy cảm và dễ bùng nổ trọng số nhất khi kết hợp nhiều nguồn đặc trưng trên tập huấn luyện chỉ có 251 người bệnh. Sự đánh đổi chỉ 0.0066 điểm C-index giữa Full CP và Full Outer-Product là một mức chi phí hoàn toàn xứng đáng để đổi lấy sự ổn định cấu trúc tại tầng tương tác.

Thứ hai, \emph{gia tăng số lượng cặp tương tác không đồng nghĩa với việc nâng cao hiệu năng dự báo}. Trong toàn bộ bảy cấu trúc CP có tương tác cặp được khảo sát, kiến trúc RC-Free chỉ duy trì hai cặp tương tác quanh ảnh MRI (MR và MC) với bốn nhánh phần dư nhưng lại xác lập C-index xác thực trung bình cao nhất (\textbf{0.8347}). Thành tích này không chỉ vượt trội hơn cấu trúc đầy đủ Full CP (0.8278) mà còn tiết kiệm được 15460 tham số; thậm chí cấu trúc đơn giản một cặp RC-only cũng đạt 0.8319, cao hơn Full CP. Nguyên nhân căn bản là do cả ba nguồn dữ liệu đầu vào vốn dĩ đã bảo tồn các đường truyền biểu diễn gốc thông qua các kết nối phần dư có cổng; vì vậy, việc cưỡng ép duy trì đồng thời cả ba cặp tương tác sẽ dẫn đến sự dư thừa thông tin và làm gia tăng phương sai ước lượng trên mẫu nhỏ. Khi đặt ảnh chụp MRI 3D làm cầu nối trung tâm liên kết sang đặc trưng kết cấu radiomics và điểm số lâm sàng, kiến trúc RC-Free đã chứng minh là giải pháp tinh gọn, cân đối và đạt hiệu quả tối ưu nhất.

Thứ ba, \emph{ý nghĩa thực tiễn của việc phân tầng nguy cơ luôn gắn liền với quy mô mẫu và số lượng biến cố tích lũy}. Phân tích sống còn Kaplan--Meier cho thấy điểm log-nguy cơ của mô hình RC-Free phân tách rất sắc nét hai nhóm bệnh nhân có diễn biến khác biệt: trên tập xác thực đạt mức ý nghĩa log-rank $p=0.0134$ (ghi nhận 12/26 ca biến cố ở nhóm nguy cơ cao so với chỉ 4/28 ca ở nhóm nguy cơ thấp) và trên tập kiểm tra hậu nghiệm độc lập đạt $p<0.001$. Tuy nhiên, khi đánh giá sâu theo từng thời hạn lâm sàng, các giá trị AUC xác thực rất cao tại mốc 2 năm (0.9805) và 3 năm (0.9788) thực chất chỉ được xây dựng trên cơ sở \textbf{7 ca biến cố tích lũy}. Đến mốc 5 năm, AUC giảm xuống 0.7889 kèm theo độ lệch chuẩn mở rộng hơn khi số lượng đối tượng còn được theo dõi thực tế sau 5 năm giảm xuống còn 19 người (Bảng~\ref{tab:horizon-support}).

\section{Giới hạn nghiên cứu}
\label{sec:study-limitations}
Mặc dù đạt được những kết quả khả quan, khóa luận nhìn nhận bốn giới hạn phương pháp luận cần lưu ý sau:
\begin{enumerate}[itemsep=0.15em,topsep=0.25em]
\item \textbf{Dữ liệu mới dừng ở bộ dữ liệu chuẩn quốc tế ADNI, chưa được kiểm chứng tại Việt Nam.} Toàn bộ quá trình huấn luyện, tinh chỉnh và kiểm thử mô hình hiện mới chỉ thực hiện trên cơ sở dữ liệu quốc tế ADNI với quy mô tập kiểm chứng còn khiêm tốn (54 người ở tập xác thực và 54 người ở tập kiểm tra hậu nghiệm). Mô hình chưa được thu thập và đánh giá trên dữ liệu bệnh nhân thực tế từ các bệnh viện lớn tại Việt Nam, nơi hệ thống máy chụp MRI, quy trình thăm khám lâm sàng cũng như đặc điểm nhân khẩu học có thể có nhiều khác biệt.
\item \textbf{Thiết kế mốc thời gian mang tính hồi cứu.} Thời gian sống còn $t_i$ được tính bắt đầu từ ngày chụp MRI lần thứ ba $p_i$, song quy tắc lọc mẫu lại đòi hỏi người bệnh phải duy trì ổn định đến hết mốc 18 tháng lịch $L_i$ (do đó tiềm ẩn nguy cơ sai lệch thời gian bất tử nếu suy diễn như một công cụ dự báo tiến cứu ngay tại ngày chụp $p_i$ \cite{suissa2008immortal}). Thêm vào đó, cửa sổ liên kết dữ liệu $\pm90$ ngày khiến cho 57 trường hợp có điểm MMSE hoặc CDR gắn với lần chụp cuối thực tế lại được đo sau ngày $p_i$ (Mục~\ref{sec:data-split}), làm cho các kết quả hiện tại phù hợp nhất cho bài toán đánh giá hồi cứu trên nhóm ổn định.
\item \textbf{Các biến đầu vào và tiêu chí xác định biến cố cùng dựa trên thang đo MMSE và CDR.} Mặc dù hai giai đoạn thu thập dữ liệu đã được phân tách rõ ràng trước và sau mốc 18 tháng, việc cùng căn cứ vào ngưỡng điểm của MMSE và CDR để xác định biến cố kết cục vẫn tạo ra sự phụ thuộc nhất định vào hai thang đo này; nghiên cứu chưa có điều kiện đối chiếu với chẩn đoán chuyên khoa độc lập hoặc các xét nghiệm sinh học thần kinh chuyên sâu hơn.
\item \textbf{Đầu ra của mô hình mới dừng lại ở việc xếp hạng nguy cơ tương đối.} Chỉ số log-nguy cơ $\eta_i$ từ mô hình Cox chỉ phản ánh thứ tự rủi ro tiến triển sớm hay muộn giữa các cá thể, song chưa quy đổi trực tiếp thành xác suất tiến triển cụ thể theo tỷ lệ phần trăm cho từng người bệnh tại các mốc thời gian mong muốn, đồng thời chưa tiến hành đo đạc thời gian đáp ứng tính toán thực tế trên các thiết bị tại cơ sở y tế.
\end{enumerate}

\section{Hướng phát triển}
Xuất phát từ các giới hạn đã được chỉ rõ ở trên, tôi đề xuất ba định hướng nghiên cứu mở rộng mang tính thực tiễn cao:
\begin{enumerate}[itemsep=0.15em,topsep=0.25em]
\item \emph{Thu thập và kiểm chứng trên dữ liệu bệnh viện thực tế tại Việt Nam}: Chủ động phối hợp với các bệnh viện lớn trong nước để thu thập dữ liệu theo dõi dọc của bệnh nhân Việt Nam, tiến hành đánh giá độ khái quát hóa ngoại kiểm và hiệu chỉnh mô hình nhằm thích ứng tốt nhất với bối cảnh lâm sàng trong nước.
\item \emph{Hoàn thiện thiết kế dự báo chuẩn tắc theo mốc thời gian}: Thiết lập giới hạn nghiêm ngặt nhằm bảo đảm toàn bộ dữ liệu lâm sàng đầu vào được đo trước ngày chụp MRI cuối $p_i$, hoặc áp dụng triệt để khung thiết kế dự báo tại mốc cố định $L_i$ (landmarking) \cite{vanhouwelingen2007dynamic} để loại bỏ hoàn toàn độ lệch thời gian giữa đầu vào và biến cố kết cục.
\item \emph{Cá thể hóa xác suất tiến triển và nâng cao tính giải thích}: Ước lượng bổ sung hàm nguy cơ nền theo phương pháp Breslow \cite{kvamme2019time} nhằm chuyển đổi điểm nguy cơ thành xác suất tiến triển chính xác theo từng năm cho mỗi bệnh nhân; đồng thời phân tích sâu các trọng số chú ý attention và cơ chế cổng điều tiết để cung cấp các giải thích trực quan, hỗ trợ đắc lực cho các bác sĩ trong quá trình hội chẩn.
\end{enumerate}

\section{Kết luận}
Bài toán tiên lượng nguy cơ tiến triển ở người mắc suy giảm nhận thức nhẹ (MCI) đòi hỏi không chỉ khả năng theo dõi sát sao tiến trình bệnh lý qua các lần khám lặp lại theo thời gian, mà còn cần sự kết hợp hài hòa giữa hình thái cấu trúc não bộ, các chỉ số định lượng toán học và đánh giá nhận thức lâm sàng. Dựa trên tập dữ liệu hồi cứu gồm 359 người tham gia từ sáng kiến ADNI qua ba mốc thăm khám dọc, khóa luận đã xây dựng thành công một quy trình học sâu phân tích sống còn đa phương thức hoàn chỉnh, tích hợp chặt chẽ ảnh thể tích MRI 3D, đặc trưng radiomics toàn não và các thang điểm lâm sàng MMSE/CDR. Toàn bộ chuỗi thực nghiệm đối chứng đã khẳng định vững chắc tính đúng đắn của từng quyết định thiết kế: chuỗi MRI dọc với mạng T-LSTM và cơ chế attention vượt trội rõ rệt so với ảnh đơn thời điểm; việc tích hợp đầy đủ cả ba nguồn dữ liệu nâng cao đáng kể năng lực phân biệt rủi ro; và cơ chế mô hình hóa tương tác song tuyến theo cặp khai thác các mối liên kết chéo hiệu quả hơn hẳn phương thức ghép nối thông thường.

Đặc biệt, kiến trúc đề xuất \textbf{RC-Free} đã kết hợp nhuần nhuyễn hai giải pháp tinh giản then chốt: kỹ thuật phân rã tensor CP hạng thấp giúp cắt giảm tới 98.38\% số lượng tham số ở ba lớp song tuyến đầu tiên, và việc loại bỏ cặp tương tác trực tiếp RC giúp tiết kiệm thêm 15460 tham số so với cấu trúc Full CP mà vẫn bảo toàn trọn vẹn cả ba nguồn dữ liệu đầu vào. Trên tập xác thực, RC-Free xác lập chỉ số C-index trung bình cao nhất trong toàn bộ bảy cấu hình khảo sát ($0.8347\pm0.0237$) và chứng minh năng lực phân tầng sống còn sắc nét, có ý nghĩa thống kê cao trên đường cong Kaplan--Meier ($p=0.0134$). Những kết quả đạt được khẳng định rằng việc kết hợp giữa tham số hóa hạng thấp và chọn lọc cấu trúc tương tác hợp lý là một hướng đi tinh gọn, chuẩn xác và đầy triển vọng cho các bài toán phân tích sống còn đa phương thức trong y tế.
